% ============================================
% Results
% ============================================
\section{Results}
\label{sec:results}

% --------------------------------------------


\subsection{Community-level selection is prevalent in microbial coalescence}
\label{sec:results:dominance_prevalent}

Our central question is whether selection acts primarily on whole communities as cohesive units during community coalescence. To address this, we asked how often coalescence outcomes are represented primarily by one parental community versus blending of both or forming a novel community. Consider a coalescence event where communities A and B merge to produce community C (\figref{fig:fig1}A). We represent each community by its normalized abundance vector and quantify the similarity between the coalesced community C and each parental community (\figref{fig:fig1}B):
\begin{equation}
\text{Sim}(C, A) = \vec{x}_C \cdot \vec{x}_A, \quad \text{Sim}(C, B) = \vec{x}_C \cdot \vec{x}_B
\label{eq:similarity}
\end{equation}
where $\vec{x}_A$, $\vec{x}_B$ are normalized abundance vectors of parental communities A and B, and $\vec{x}_C$ is that of the coalesced community. These two similarity scores yield a two-dimensional similarity map in which the location of C reflects how strongly it resembles each parental community. We partition this map into three coalescence outcome classes (see Methods): Dominance, where C closely resembles one parental community but not the other; Mixture, where C similarly resembles both parental communities; and Restructuring, where C diverges from both parental communities into a novel configuration. Crucially, Dominance reflects community-level selection: one parental community displaces the other while retaining its internal compositional structure, indicating that its species persist collectively.

To quantify the occurrence of outcome types experimentally, we assembled random in-vitro communities and subjected them to pairwise coalescence (\figref{fig:fig1}C). We first curated a strain library of 54 bacterial isolates collected from diverse environments (soil, tree surface, and flower stamen). The library is phylogenetically broad, spanning 29 families across three phyla: Proteobacteria, Firmicutes, and Bacteroidota (Extended Data Fig.~1; Supplementary Fig.~1). From this library, we assembled 30 parental communities with varying initial richness (6, 12, or 24 species; see Methods) and stabilized them for 7 days under daily growth--dilution cycles ($\times$30 every 24~h) (\figref{fig:fig1}C). Our initial experiments were done in Base medium (1~g~L$^{-1}$ yeast extract, 1~g~L$^{-1}$ soytone, 10~mM sodium phosphate, 0.1~mM CaCl$_2$, 2~mM MgCl$_2$, 4~mg~L$^{-1}$ NiSO$_4$, 50~mg~L$^{-1}$ MnCl$_2$, 5~g~L$^{-1}$ glucose and 4~g~L$^{-1}$ urea; pH~6.5), a buffered complex medium that we have previously determined has moderate interaction strength and supports high species coexistence (Supplementary Information, mean species survival ratio = $74 \pm 2$\%). We then performed 83 pairwise coalescence events by mixing stabilized parental communities at equal volume and restabilizing for an additional 7 days. Community compositions of parental and coalesced communities were measured at the end of stabilization by 16S rRNA amplicon sequencing (Amplicon Sequence Variants, ASVs), and we also recorded the optical density (OD$_{600}$) and pH of the communities to contextualize assembly and post-coalescence dynamics.

Representative time series illustrate the spectrum of post-coalescence dynamics (\figref{fig:fig1}D). In outcomes classified as Dominance, one parental lineage rapidly displaces the other after mixing and the coalesced community converges to a composition closely matching that parental community, while largely preserving its internal relative-abundance structure over serial dilution cycles. In the Restructuring case, the merged community converges toward a novel state distinct from both parental communities. Among 3 representative time-series trajectories, we did not observe Mixture cases in which both parental communities remained comparably represented after coalescence (see Supplementary Fig.~26).

Projecting all outcomes into the similarity space based on stabilized compositions (\figref{fig:fig1}E) revealed that Dominance is the most frequent empirical outcome. Of the 83 coalescence events, 54 were classified as Dominance, 26 as Restructuring, and 3 as Mixture (\figref{fig:fig1}E). This pattern of Dominance as the most frequent outcome was robust across variants of similarity metrics (Extended Data Fig.~2). One possibility is that the observed frequency of Dominance results from skewed species abundance distributions rather than correlated selection among species from the same parental community. To rule out this possibility, we compared experimental outcomes against two null models that assume no correlation in species selection: (1) abundance-weighted random selection and (2) shuffled abundance (see Supplementary Information). The experimentally observed asymmetry significantly exceeded both null expectations (Extended Data Fig.~3), supporting that Dominance reflects correlated species selection within parental communities.

% Figure 1
\begin{figure}[H]
\centering
\includegraphics[width=\textwidth]{figure_main_1_experimental_overview.pdf}
\caption{\textbf{Fig.~1. Coalescence of synthetic microbial communities frequently yields Dominance.} \textbf{a}, Schematic of coalescence: parental communities A and B are mixed to produce C. Outcomes are classified as Dominance (one parent wins), Mixture (both persist), or Restructuring (novel state). \textbf{b}, We use similarity to quantify compositional resemblance between communities and define a two-dimensional similarity map for classifying outcomes. \textbf{c},  Experimental workflow: 30 parental communities (initial richness 6, 12, or 24) were assembled from 54 bacterial isolates, stabilized, mixed pairwise ($n = 83$), and restabilized. \textbf{d}, Representative time courses showing Dominance (left, center) and Restructuring (right). \textbf{e}, Dominance is the most frequent outcome in Base medium. Of 83 coalescence events, 54 (65\%) were classified as Dominance, 26 (31\%) as Restructuring, and 3 (4\%) as Mixture. Different symbols indicate initial richness (circles: 6; squares: 12; triangles: 24 species).}
\label{fig:fig1}
\end{figure}

% --------------------------------------------
\subsection{Theoretical model with random interactions reproduces community-level selection}
\label{sec:results:glv_model}
To gain insight into why Dominance is the prevalent outcome, we introduced a generalized Lotka--Volterra (gLV) model \citep{May1972, Grilli2017} that mirrors the experimental protocol (\figref{fig:fig2}A). In this classic model, species grow logistically and compete pairwise:
\begin{equation}
\frac{dn_i}{dt} = n_i \left( r_i - \sum_j \alpha_{ij} n_j \right)
\label{eq:glv}
\end{equation}
where $n_i$ is abundance, $r_i$ is growth rate, and $\alpha_{ij}$ is the competition coefficient between species $i$ and $j$ (with self-interaction $\alpha_{ii} = 1$). We fixed $r_i = 1$ and drew off-diagonal competition coefficients from a uniform distribution $\mathbb{U}(0, 2\mu)$ with mean $\mu$, which controls the average competition strength: higher $\mu$ means stronger interspecies competition \citep{Hu2022,Hu2025}. From a pool of 54 species, we generated two parental communities of 12 species each with no shared species, allowed each to reach equilibrium, then mixed them equally to simulate coalescence (\figref{fig:fig2}A; see Methods for details). Post-coalescence compositions were analyzed using the same similarity metrics as in the experiments.

We simulated 1,200 random coalescence events at interaction strength $\mu = 0.6$ and analyzed the similarity of each coalescence outcome to the two parental communities. At this interaction strength, the random-interaction model quantitatively reproduces the high frequency of Dominance observed in the experiments. Outcomes concentrate near the axes rather than the diagonal, indicating asymmetric dominance by one parental community (\figref{fig:fig2}B). Quantitatively, Dominance accounts for 61\% of all outcomes, far more frequent than Restructuring (26\%) or Mixture (13\%; \figref{fig:fig2}B, right). This observation suggests that interspecies interactions alone, a minimal and generic condition, are sufficient to reproduce the high frequency of Dominance observed in coalescence.

To understand how Dominance emerges in the model and whether it reflects community-level selection, we examined the role of the assembly process. During assembly, competitive exclusion filters out species with strong competitive interactions, leaving communities with reduced mean interaction strength compared to the initial pool (paired $t$-test, $P <0.001$; \figref{fig:fig2}C). Because surviving species compete weakly with each other but face stronger competition from outsiders, their fates become coupled during coalescence: they tend to persist or go extinct together. We quantified this effect using pairwise selection correlation—the degree to which species pairs share the same survival outcome (both persist or both go extinct) across coalescence events (Supplementary Information; \figref{fig:fig2}D). Species from the same parental community showed strongly positive correlations, while cross-community pairs showed negative correlations—a pattern observed in both simulations and experiments (\figref{fig:fig2}D). This confirms that member species within a parental community, having undergone assembly together, have coupled fates and are co-selected during coalescence. Together, the prevalence of Dominance combined with positive within-community selection correlation provides evidence for community-level selection in this system.


% Figure 2
\begin{figure}[H]
\centering
\includegraphics[width=\textwidth]{figure_main_2_simulation_framework.pdf}
\caption{\textbf{Fig.~2. Generalized Lotka--Volterra model of community coalescence.} \textbf{a}, Simulation workflow: species grow according to the generalized Lotka--Volterra (gLV) equations, where interaction coefficients $\alpha_{ij}$ are drawn from a uniform distribution $\mathbb{U}(0, 2\mu)$ with mean $\mu$ controlling mean interaction strength. Two equilibrated parental communities are mixed and simulated to steady state. \textbf{b}, Outcome density map ($\mu = 0.6$, $n = 1{,}200$ simulations). The model reproduces the high frequency of Dominance (61\%) observed experimentally. \textbf{c}, Community assembly significantly reduces the mean interaction strength among surviving species (paired $t$-test, $p < 0.001$). \textbf{d}, Pairwise selection correlation. Species pairs from the same parental community show positive correlation (co-survival), while cross-community pairs show negative correlation, indicating community-level selection. This pattern holds for both simulations (left) and experiments (right). Error bars, s.e.m.}
\label{fig:fig2}
\end{figure}


% --------------------------------------------
\subsection{Interaction strength controls coalescence outcome type and degree of community-level selection}
\label{sec:results:interaction_strength}

Having established that random interspecies interactions can recapitulate Dominance with correlated pairwise selection, we next investigated how interaction strength alters coalescence outcomes. We simulated 1,200 coalescence events at each of three representative interaction strengths ($\mu = 0.3, 0.6, 0.8$) and mapped outcomes into the similarity space (\figref{fig:fig3}A). At weak interactions ($\mu = 0.3$), outcomes clustered in the interior of the map, indicating Mixture; as $\mu$ increased, outcomes shifted toward the axes, indicating Dominance. We quantified this shift using the Parental Dominance Index (PDI), which captures the relative contribution of each parental community (0: parental community B, 0.5: equal, 1: parental community A; see Methods). PDI distributions shifted from unimodal near 0.5 at $\mu = 0.3$ to bimodal at higher $\mu$ (\figref{fig:fig3}A), reflecting a transition in which Dominance becomes increasingly frequent, indicating community-level selection.

Across interaction strengths from $\mu = 0$ to $1.2$, we observed a systematic shift in coalescence outcomes (\figref{fig:fig3}B), with Mixture predominating at weak interactions and Dominance at strong interactions. Restructuring emerged at moderate to high interaction strengths. These patterns were robust to variation in carrying capacities, interaction distributions, similarity metrics, and community size (Supplementary Figs.~2–4; Extended Data Fig.~4). Notably, the frequency of Dominance remained relatively stable across parental communities ranging from 4 to 48 species, spanning our experimental species pool range. To determine whether this transition from Mixture to Dominance corresponds to the degree of co-selection, we examined pairwise selection correlation across interaction strengths. Notably, within-community pairwise selection correlation emerged only at high interaction strengths, indicating that species fates become coupled only when interactions are sufficiently strong (Extended Data Fig.~5). These findings indicate that interaction strength simultaneously determines both coalescence outcome type and the level at which selection operates.

% Figure 3
\begin{figure}[H]
\centering
\includegraphics[width=\textwidth]{figure_main_3_interaction_strength.pdf}
\caption{\textbf{Fig.~3. Interaction strength controls the transition between coalescence outcome types in simulation.} \textbf{a}, Simulated coalescence outcome maps at three representative interaction strengths: $\mu = 0.3$, $0.6$, $0.8$; $n = 1{,}200$ simulations per condition. Outcomes shift from central Mixtures at weak interactions ($\mu = 0.3$) to axis-aligned Dominance at strong interactions ($\mu = 0.8$). Histograms show the Parental Dominance Index (PDI) transitioning from unimodal to bimodal. \textbf{b}, Outcome fractions across interaction strengths ($\mu = 0$--$1.2$). Mixture predominates at low $\mu$, while Dominance becomes prevalent as $\mu$ increases. Restructuring peaks at intermediate strengths. Error bars, s.e.m.}
\label{fig:fig3}
\end{figure}


% --------------------------------------------
\subsection{Nutrient-dependent interaction strength in experiments recapitulates model predictions}
\label{sec:results:nutrient_experiment}

Following prior work showing that nutrient concentration intensifies microbial competition \citep{Hu2022, Ratzke2020, Hu2025}, we conducted additional coalescence experiments by removing or augmenting glucose and urea in the Base medium used in \figref{fig:fig1} (Methods). This yielded two additional media conditions (\figref{fig:fig4}A): Nutr$-$ (no added glucose/urea) and Nutr$+$ (high supplementation). Higher nutrient concentration amplifies consumer--resource feedbacks and intensifies environmental modification, thereby strengthening interspecies interactions \citep{Ratzke2020}. To empirically validate this effect, we measured failed invasion frequency using pairwise invasion assays among the 12 most abundant isolates (95:5 initial frequency; Methods, Supplementary Figs.~5–7). The fraction of failed invasions increased monotonically with nutrient supply (Nutr$-$: $2 \pm 1\%$; Base: $33 \pm 4\%$; Nutr$+$: $48 \pm 4\%$; mean $\pm$ s.e.m.; \figref{fig:fig4}B). Assuming the uniform distribution used in the model, calibrating these values against gLV simulations yielded approximate mean interaction strengths of $\mu \approx 0.5$ for Nutr$-$, $\mu \approx 0.7$ for Base, and $\mu \approx 0.9$ for Nutr$+$ (see Supplementary Methods).

Given that nutrient concentration modulates interaction strength, we performed coalescence experiments in Nutr$-$ and Nutr$+$ media using the same parental community library to examine how interaction strength affects coalescence outcomes. The distribution of outcomes in similarity space shifted systematically with nutrient concentration (\figref{fig:fig4}C), which we quantified by the fraction of each outcome type (\figref{fig:fig4}D). In Nutr$-$ ($n = 90$), where interactions are weakest, Mixture was the most frequent outcome (53\%), and Dominance was substantially reduced to 39\% compared to the Base medium (Dominance: 65\%, Mixture: 4\%). In Nutr$+$ ($n = 90$), Dominance further increased to 76\%. Pairwise selection correlation analysis confirmed that this shift corresponds to a transition in the level of selection (Extended Data Fig.~6). In Base and Nutr+ media, within-community species pairs showed significantly higher selection correlation than cross-community pairs ($P < 0.001$), consistent with community-level selection. In Nutr$-$, no such correlation was observed, indicating species-level dynamics. These results experimentally validate the theoretical prediction: weaker interactions yield Mixture with uncorrelated species fates, while stronger interactions yield Dominance with community-level selection.

% Figure 4
\begin{figure}[H]
\centering
\includegraphics[width=\textwidth]{figure_main_4_nutrient_perturbation.pdf}
\caption{\textbf{Fig.~4. Nutrient concentration modulates interaction strength and validates model predictions.} Experimental manipulation of nutrient concentration confirms that stronger interactions shift coalescence outcomes from Mixture toward Dominance. \textbf{a}, Schematic of nutrient manipulation: we varied glucose and urea concentrations to create three media conditions with increasing interaction strength (Nutr$-$, Base, Nutr+). \textbf{b}, Pairwise invasion assays confirm that nutrient concentration modulates interaction strength. Failed invasion frequency among the 12 most abundant isolates (95:5 initial frequency, $n = 132$ assays per medium) increases monotonically with nutrient concentration (Nutr$-$: $2 \pm 1\%$, Base: $33 \pm 4\%$, Nutr$+$: $48 \pm 4\%$; mean $\pm$ s.e.m.). \textbf{c}, Coalescence outcome distributions shift systematically with nutrient concentration. Scatter plots show outcomes in similarity space for each medium (Nutr$-$: $n = 90$, magenta; Base: $n = 83$, brown; Nutr+: $n = 90$, orange); different symbols indicate initial richness. Histograms show PDI distributions. \textbf{d}, The fraction of Dominance increases with nutrient concentration (39\% in Nutr$-$, 65\% in Base, 76\% in Nutr+) while Mixture declines from 53\% to 4\% to 6\% (chi-square test for trend, $p < 0.001$), consistent with model predictions. Error bars, s.e.m.}
\label{fig:fig4}
\end{figure}


% --------------------------------------------
\subsection{Dominant community predictability reveals two mechanistic regimes}
\label{sec:results:predictability}

We observed Dominance in both Base and Nutr$+$ media, consistent with community-level selection as confirmed by pairwise selection correlation (Extended Data Fig.~6). We next asked whether we can predict which parental community will win during coalescence. Prior work has suggested that species-level competitive outcomes may correlate with Dominance direction \citep{Rillig2015, Castledine2020}. Inspired by these observations, we tested whether pairwise competition between dominant species \citep{Hu2022, Ratzke2020} predicts which community wins (\figref{fig:fig5}A).

The relative abundance of dominant species in stabilized parental communities increased with nutrient concentration ($44 \pm 2\%$ in Nutr$-$, $51 \pm 5\%$ in Base, $67 \pm 4\%$ in Nutr$+$; mean $\pm$ s.e.m., \figref{fig:fig5}B), suggesting that dominant species exert greater influence under higher nutrient concentration. We therefore tested whether pairwise competition between dominant species predicts the coalescence outcome by analyzing the linear relationship between dominant-species competitive success (from invasion assays) and the PDI of coalescence outcomes (\figref{fig:fig5}C). Predictive power varied markedly across media: in Nutr$-$, where Mixture predominates and pairwise selection correlation is weak, pairwise competition showed no predictive power ($R^2 = 0.00$); in Base medium, predictive power was weak ($R^2 = 0.11$), consistent with collective multi-species dynamics shaping outcomes; in Nutr$+$, pairwise competition was more predictive ($R^2 = 0.49$), indicating that the dominant species contributes substantially to determining which community wins. These results suggest that community-level selection spans a mechanistic continuum: at one end, an emergent regime (Base medium) where multi-species dynamics collectively shape the outcome and no single species trait is predictive, and at the other, a top-down regime (Nutr$+$ medium) where a few dominant taxa determine the winner. 

We explored the mechanistic basis of the top-down regime by examining environmental modification through pH. In our system, dominant species are often strong pH modifiers that either acidify or alkalinize the medium (Supplementary Fig.~8). In both Base and Nutr$+$ media, when these pH-modifying species become dominant within a community, they determine the community's overall pH; communities dominated by acidifiers become acidic, while those dominated by alkalizers become alkaline (Supplementary Fig.~9). In coalescence events between acidic (pH $<$ 6.5) and alkaline (pH $>$ 7.5) parental communities, the acidic community won in only 56\% of cases in Base medium ($n = 41$), but won in 91\% of cases in Nutr$+$ medium ($n = 32$, Fisher's exact test $p < 0.0001$; Extended Data Fig.~8), consistent with the hypothesis that high nutrients amplify metabolic activity, thereby intensifying pH modification and excluding pH-sensitive species\citep{Ratzke2018, Ratzke2020}. Thus, in Nutr$+$ medium, the dominant species determines community pH, and community pH predicts coalescence outcome, providing a mechanistic explanation for the top-down regime.


% Figure 5
\begin{figure}[H]
\centering
\includegraphics[width=\textwidth]{figure_main_5_predictability.pdf}
\caption{\textbf{Fig.~5. Predictability of Dominance direction reveals two mechanistic regimes.} The direction of Dominance (which parental community wins) becomes increasingly predictable from dominant species competition across media conditions, distinguishing an emergent regime (Base) from a top-down regime (Nutr+). \textbf{a}, Schematic: can pairwise competition between dominant species predict which community wins? \textbf{b}, Relative abundance of dominant species increases with nutrient concentration (Nutr$-$: $44 \pm 2\%$; Base: $51 \pm 5\%$; Nutr+: $67 \pm 4\%$; mean $\pm$ s.e.m.). \textbf{c}, Dominant species competitive success versus PDI of coalescence outcomes; gray shading indicates Dominance. Predictive power: $R^2 = 0.00$ in Nutr$-$, $R^2 = 0.11$ ($p = 0.03$) in Base (emergent regime), and $R^2 = 0.49$ ($p < 0.001$) in Nutr+ (top-down regime). Linear regression with 95\% confidence intervals shown.}
\label{fig:fig5}
\end{figure}


% --------------------------------------------
\subsection{Interaction-dependent coalescence outcomes generalize to natural communities}
\label{sec:results:natural_communities}

The synthetic communities described above were assembled from individually isolated bacterial strains, allowing precise control over initial composition but raising the question of whether these patterns generalize to more ecologically realistic settings. To address this, we performed coalescence experiments using communities derived from natural environmental samples, which harbor complex species assemblages shaped by evolutionary and ecological processes in their natural habitats.

We collected six environmental samples from diverse microhabitats (soil, compost, and decomposing organic matter) and established bacterial communities through seven rounds of serial growth-dilution in laboratory media across all three nutrient conditions (Nutr$-$, Base, and Nutr+; \figref{fig:fig6}A). After stabilization, these natural sample-derived communities exhibited higher ASV richness than synthetic communities (mean of $13.7 \pm 7.2$ ASVs above 0.1\% threshold, compared to $9.8 \pm 4.8$ in synthetic communities) and low ASV overlap among communities from different samples (Supplementary Figs.~22–25).

We performed 15 pairwise coalescence events with two biological replicates each ($n = 30$ per condition) across all three nutrient conditions after stabilization and analyzed outcomes using the same similarity framework applied to synthetic communities. Natural sample-derived communities showed similar patterns to synthetic communities: outcomes clustered along the axes in similarity space (\figref{fig:fig6}B,C). The fraction of Dominance outcomes increased with nutrient concentration (37\% in Nutr$-$, 70\% in Base, 77\% in Nutr+), mirroring the interaction-strength dependence observed in synthetic communities. Natural communities showed higher Restructuring fractions, possibly reflecting greater taxonomic diversity and more complex interaction networks. These results demonstrate that the interaction-dependent shift toward community-level selection is robust and generalizes beyond precisely controlled synthetic consortia to naturally-evolved microbial assemblages.

% Figure 6
\begin{figure}[H]
\centering
\includegraphics[width=\textwidth]{figures/figure_main_6_natural_communities.pdf}
\caption{\textbf{Fig.~6. Dominance generalizes to natural sample-derived communities.} Natural environmental communities show qualitatively similar coalescence patterns to synthetic communities, with Dominance frequency increasing under higher nutrient concentration. \textbf{a}, Experimental workflow follows Fig.~1c, but uses natural sample-derived communities from six environmental samples including soil, compost, and decomposing organic matter. \textbf{b}, Coalescence outcome distributions in similarity space for natural communities across three media conditions ($n = 30$ coalescence events per condition, 15 unique pairs $\times$ 2 replicates). Similar to synthetic communities, outcomes cluster near the diagonal under weak interactions and shift toward the axes under strong interactions. Histograms below show PDI distributions. \textbf{c}, The fraction of Dominance outcomes increases with nutrient concentration in natural communities (37\% in Nutr$-$, 70\% in Base, 77\% in Nutr+), consistent with patterns observed in synthetic communities. Error bars, s.e.m.}
\label{fig:fig6}
\end{figure}
